\tocsection{\textbf{9/26/26 Lecture 9:} Rao-Blackwellization}

Recall that with estimation, we will be using $\delta(X)$, an estimate ($\delta$ is called the estimator), to guess $g(\theta)$, an estimand. To test how good a given estimate is, we compute the loss $L(\theta, \delta(X))$. Then, to test how good an estimator is, we compute the risk:
\[ R(\theta, \delta) = \bb{E}_{\theta}[L(\theta, \delta(X))] \]
There is no $\delta$ that minimizes risk everywhere, so we need to find a way to judge which estimators are better than others.\\

One way to do it is to summarize the risk with a value, whether it be the mean risk, or maximum possible risk.\\

Another way to do it is by restricting to a smaller class of estimators. We will first see what happens when we restrict to unbiased estimators. Recall that the bias of an estimator, $\tt{Bias}(\theta, \delta)$ is defined to be
\[ \tt{Bias}(\theta, \delta) = \bb{E}_{\theta}[\delta(X)] - g(\theta) \]
we say an estimator is \emph{unbiased} if $\tt{Bias}(\theta, \delta) = 0$, so the average value of $\delta(X)$ is always $g(\theta)$.

\subsection*{Convex Loss}

\begin{definition}[\textbf{Definition 9.1: }Convex Functions]
    A function $f(x)$ is \emph{convex} if, for all $x_1, x_2$ and $\gamma \in (0, 1)$, we have that
    \[ f(\gamma x_1 + (1 - \gamma)x_2) \leq \gamma f(x_1) + (1 - \gamma) f(x_2) \]
    The function $f$ is \emph{strictly convex} if equality never holds.
\end{definition}
Here is what a convex (and in this case, strictly convex) function looks like visually. These secant lines always lie above the function.
\fig{0.6}{Lecture9FIg1}{}

\begin{theorem}[\textbf{Theorem 9.2: }Jensen's Inequality]
    If $f$ is convex, then $f(\bb{E}(X)) \leq \bb{E}f(X)$. If $f$ is strictly convex, then $f(\bb{E}(X)) < \bb{E}f(X)$ unless $X$ is almost surely constant. 
\end{theorem}

This is useful since if we have a convex loss function, then we can find a lower bound on the risk by just taking the expectation of $X$ and then plugging into the loss.\\

Note: We say that $L(\theta, d)$ is a convex loss function if it is convex in $d$ (as we typically fix $\theta$).\\

\begin{example}{\textbf{Example}}
    Consider the squared error loss $L(\theta, d) = (g(\theta) - d)^2$, and the mean square error $\tt{MSE}(\theta; \delta) = \bb{E}_{\theta}[(g(\theta) - \delta(X))^2]$. Recall that we can write
    \[ \tt{MSE}(\theta; \delta) = \tt{Bias}(\theta, \delta)^2 + \tt{Var}_{\theta}(\delta(X)) \]
    Now, if $\delta$ is unbiased, then
    \[ \tt{MSE}(\theta; \delta) = \tt{Var}_{\theta}(\delta(X)) \]
    so the best unbiased estimator is the one that gives the least variance.
\end{example}

\begin{theorem}[\textbf{Theorem 9.3:} Rao-Blackwell Theorem]
    Assume that $T(X)$ is sufficient, and $\delta(X)$ is any estimator. Let 
    \[ \overline{\delta}(T(X)) = \bb{E}[\delta(X) \given T(X)] \]
    If $L(\theta, d)$ is convex, then $R(\theta, \overline{\delta}) \leq R(\theta, \delta)$ for all $\theta$. Then, if $L(\theta, d)$ is strictly convex, then $R(\theta, \delta) < R(\theta, \delta)$ unless $\delta(X) \overset{\tt{a.s.}}{=} \overline{\delta}(T(X))$ for all $\theta$. We call $\overline{\delta}(T(x))$ is called the \emph{Rao-Blackwellization} of $\delta(x)$. 
\end{theorem}

\begin{definition}[\textbf{Definition 9.4:} U-estimability]
    An estimand $g(\theta)$ is \emph{U-estimable} if there is some estimator $\delta(X)$ for which $\tt{Bias}_{\theta}(\theta, \delta) = 0$ for all $\theta$
\end{definition}

\begin{example}{\textbf{Example}}
    Let $X \sim \Bern(\theta)$, $\theta \in (0, 1)$. Now, $\bb{E}_{\theta}[\delta(X)] = (1 - \theta)\delta(0) + \theta \delta(1)$, so estimators are just maps $\{0, 1\} \to \R$, and there is no way to have this be zero for all $\theta$, so $\theta$ is not U-estimable. 
\end{example}

\begin{definition}[\textbf{Definition 9.5:} Uniform Minimum Variance Unbiased (UMVU)]
    We say that $\delta(X)$ is \emph{uniform minimum variance unbiased (UMVU)} if for any unbiased $\tilde{\delta}(X)$
    \[ \Var_{\theta}(\delta(X)) \leq \Var_{\theta}(\tilde{\delta}(X)) \]
    for all $\theta$.
\end{definition}

\begin{theorem}[\textbf{Theorem 9.6:} Uniqueness of Rao-Blackwellization]
    Assume that $T(X)$ is complete sufficient, $g(\theta)$ is U-estimable and $L(\theta, d)$ is convex. Then, there is a unique unbiased estimator $\delta^{*}$ of the form $\delta^{*}(T(X))$ that uniformly minimizes risk (over all $\theta$) among all unbiased estimators. 
\end{theorem}
\begin{proof}{Proof Sketch:}
    Since there is some unbiased estimator as $g(\theta)$ is U-estimable, then we can Rao-Blackwellize it. Then, it must be unique as $T(X)$ is complete. Then, it uniformly minimizes risk over all unbiased estimators because the Rao-Blackwellization of all unbiased estimators $\delta$ is $\delta^{*}$, and this has less risk than $\delta^{*}$. 
\end{proof}
Notice that Theorem 9.6 implies that $\delta^{*}$ also is UMVU as it minimizes the MSE and is unbiased.\\

Now, let's go into some examples to show the general process for finding the UMVU estimator. We will consider the case when $X_1, \ldots, X_n \iidsim \Pois(\theta)$. Notice that $\overline{X}$ is the UMVU for $g(\theta) = \theta$ as it is an unbiased estimator which is a function of the complete sufficient statistic (recall that $T(X) = \sum x_i$ is the complete sufficient statistic)\\

Now, what about estimators for $g(\theta) = \theta^2$? Is it naively $\overline{X}^2$? Well, let's use our methods to find it. Recall that the probability mass function as a function of $T(X)$ is
\[ p_{\theta}(t) = \frac{e^{-n\theta}(n\theta)^{t}}{t!} \]
We will find the UMVU in two ways: Solving for an unbiased estimator based on $T$, and Rao-Blackwellizing any unbiased estimator.\\

\textbf{Strategy 1:} Create an Estimator as a function of $T$\\

Now, suppose that we have an unbiased estimator $\delta(T)$ as a function of the complete sufficient statistic. Then, we know that $\delta(T)$ is unbiased if, and only if
\begin{align*}
    \theta^2 &= \int_{\{0, 1, \ldots\}} \delta(t)dP_{\theta}(t)\\
    &= \sum_{t = 0}^{\infty} \delta(t)p_{\theta}(t)\\
    &= e^{-n\theta}\sum_{t = 0}^{\infty} \delta(t)\frac{n^{t}}{t!}\theta^t
\end{align*}
Thus, we see that $e^{n\theta} \theta^2 = \sum_{t = 0}^{\infty} \delta(t)\frac{n^{t}}{t!}\theta^t$. Then, we can expand the left hand side of this equation to get two power series
\[ \sum_{k = 0}^{\infty} \frac{n^{k}\theta^{k + 2}}{k!} = \sum_{t = 0}^{\infty} \delta(t)\frac{n^{t}\theta^{t}}{t!} \]
Now, change the index so the powers of $\theta$ start at the first term, which gives
\[ \sum_{t = 2}^{\infty} \frac{n^{t - 2}\theta^{t}}{(t - 2)!} = \sum_{t = 0}^{\infty} \delta(t)\frac{n^{t}}{t!}\theta^{t} \]
Matching terms, we must have that $\delta(0) = 0$, $\delta(1) = 0$, and for $t \geq 2$, we must have
\[ \delta(t) = \frac{t(t - 1)}{n^2} \]
Thus, that is our estimator!
\[ \delta(T(X)) = \frac{T(X)(T(X) - 1)}{n^2} = \overline{X}^2 - \frac{\overline{X}}{n} \]
since $\overline{X} = \frac{T(X)}{n}$.\\

\textbf{Strategy 2:} Rao-Blackwellize an Unbiased Estimator\\

Let $\delta(X) = X_1 \cdot X_2$. Since $\bb{E}[\delta(X)] = \theta^2$, this is an unbiased estimator (although a bad one). Now, let's Rao-Blackwellize. Let
\[ \delta^{*}(T) = \bb{E}[X_1X_2 \given T] \]
Recall that $X \given T \sim \tt{Multinom}(t, \frac{1}{n}1_{n})$ where $1_{n}$ is a $n$-vector of $1$'s. Thus, we see that $X_1 \given T \sim \tt{Binom}(t, \frac{1}{n})$. Thus, we see that
\[ \bb{E}[X_1 \given T] = \frac{T}{n} \quand \tt{Var}(X_1 \given T) = T \cdot \frac{1}{n} \cdot \left(1 - \frac{1}{n}\right) = \frac{T(n - 1)}{n^2} \]
Then, as $X_2 \given T, X_1 \sim \Binom(t - x_1, \frac{1}{n - 1})$, we have that
\[ \bb{E}[X_2 \given T, X_1] = \frac{T - X_1}{n - 1} \]
Thus, we have that
\begin{align*}
    \bb{E}[X_1X_2 \given T] &= \bb{E}[X_1 \cdot \bb{E}[X_2 \given T, X_1] \given T]\\
    &= \bb{E}\left[X_1 \cdot \frac{T - X_1}{n - 1} \given T\right]\\
    &= \bb{E}\left[\frac{TX_1}{n - 1} - \frac{X_1^2}{x - 1} \given T\right]\\
    &= \frac{T(T - 1)}{n^2}
\end{align*}
which is the same thing we got before!\\

\textbf{Trick Strategy 3:} Fun trick\\

There is an even easier way to do this (just by coincidence that $\overline{X}^2$ is a moment). Notice that as 
\[ \bb{E}[\overline{X}^2] = \Var(\overline{X}) + (\bb{E}[\overline{X}])^2 = \frac{\theta}{n} + \theta^2 \]
we see that
\[ \bb{E}\left[\overline{X}^2 - \frac{\overline{X}}{n}\right] = \theta^2 \]
so as this is the expectation of $\delta(T) = \frac{T^2}{n^2} - \frac{T}{n^2}$, we see it is unbiased and as it is a function of $T$, is is the UMVU. \\

Here is one last interesting example.

\begin{example}{\textbf{Example}}
    Let $X_1, \ldots, X_n \iidsim U[0, \theta]$ with $\theta > 0$. Now, recall that $T = X_{(n)}$ is complete sufficient with
    \[ p_{\theta}^{T}(X) = \frac{n}{\theta^{n}}t^{n - 1}1\{t \leq \theta\} \]
    Then, notice we have that
    \begin{align*}
        \bb{E}_{\theta}[T] &= \int_{0}^{\theta} tp_{\theta}(t) dt\\
        &= \int_{0}^{\theta} \frac{n}{\theta^{n}}t^{n}dt\\
        &= \left[\frac{nt^{n + 1}}{(n + 1)\theta^{n}}\right]_{0}^{\theta}\\
        &= \frac{n}{n + 1}\theta
    \end{align*}
    Thus, we then see that
    \[ \bb{E}_{\theta}\left[\frac{n + 1}{n}T\right] = \theta \]
    so $\delta(T) = \frac{n + 1}{n}T$ must be the UMVU as it is an unbiased estimator of $T$.

    $\qed$
\end{example}



